%%%%%%%%%%%%%%%%%%%%%%%%
%
% $Autor: Hemanth Jadiswami Prabhakaran $
% $Datum: 2026-09-25 $
% $Pfad: GitHub/MBIDA_Master_Thesis/Manual/Chapters/en/02Installation.tex $
% $Version: 1 $
%
% $Project: MBIDA Thesis - Manual $
%
%%%%%%%%%%%%%%%%%%%%%%%%


\chapter{Installation}
\label{chap:install}

There are two ways to use the dashboard (Figure~\ref{fig:installPaths}). \textbf{Cloud access} (Section~\ref{sec:instCloud}) needs nothing but a web browser and is the right choice for almost every user. \textbf{Local installation} (Section~\ref{sec:instLocal}) runs the same dashboard on your own computer; it is needed only to work offline, to check the installation independently of the cloud, or to prepare changes (Chapter~\ref{chap:maintenance}).

\begin{figure}[H]
  \centering
  \resizebox{0.95\linewidth}{!}{% Decision flow: which way of accessing the dashboard to choose.
\begin{tikzpicture}[node distance=0.75cm and 1.2cm]

  \node[man/terminal] (start) {Start};
  \node[man/decision,below=of start,text width=3.0cm] (q1) {Do you only want to\\\emph{use} the dashboard?};

  % --- cloud path (left) ---
  \node[man/boxgreen,below left=1.0cm and 0.4cm of q1,text width=4.2cm] (c1)
      {\textbf{A \; Cloud access}\\[2pt]Open the link in a browser};
  \node[man/box,below=of c1,text width=4.2cm] (c2) {Wake the app if it is asleep\\\footnotesize(first visit after a quiet period)};
  \node[man/box,below=of c2,text width=4.2cm] (c3) {Dashboard ready\\\footnotesize nothing to install};

  % --- local path (right) ---
  \node[man/boxorange,below right=1.0cm and 0.4cm of q1,text width=4.4cm] (l1)
      {\textbf{B \; Local installation}\\[2pt]Install Python ($\geq$ 3.10) and Git};
  \node[man/box,below=of l1,text width=4.4cm] (l2) {Clone the repository};
  \node[man/box,below=of l2,text width=4.4cm] (l3) {Create a virtual environment\\and install the 5 packages};
  \node[man/decision,below=of l3,text width=3.1cm] (l4) {Are the 4 result files\\present?};
  \node[man/boxred,right=0.6cm of l4,text width=2.4cm] (l4b) {Chapter~7:\\regenerate them};
  \node[man/box,below=of l4,text width=4.4cm] (l5) {\texttt{streamlit run \dots}\\\footnotesize opens \texttt{localhost:8501}};

  \node[man/terminal,below=1.0cm of $(c3)!0.5!(l5)$] (end) {Continue with Chapter 3};

  \draw[man/arrow] (start) -- (q1);
  \draw[man/arrow] (q1.west) -| node[man/label,above,pos=0.25]{yes} (c1.north);
  \draw[man/arrow] (q1.east) -| node[man/label,above,pos=0.25]{no -- offline use,\\ changes, examiner check} (l1.north);
  \draw[man/arrow] (c1) -- (c2);
  \draw[man/arrow] (c2) -- (c3);
  \draw[man/arrow] (l1) -- (l2);
  \draw[man/arrow] (l2) -- (l3);
  \draw[man/arrow] (l3) -- (l4);
  \draw[man/arrow] (l4) -- node[man/label,right]{yes} (l5);
  \draw[man/arrow] (l4) -- node[man/label,above]{no} (l4b);
  \draw[man/dashedarrow] (l4b) |- (l5.east);
  \draw[man/arrow] (c3) |- (end.west);
  \draw[man/arrow] (l5) |- (end.east);

\end{tikzpicture}
}
  \caption{Choosing between cloud access and local installation}
  \label{fig:installPaths}
\end{figure}

\section{Option A: Cloud Access (Recommended)}
\label{sec:instCloud}

\subsection{Requirements}

\begin{itemize}
  \item A current version of Chrome, Firefox, Edge or Safari with JavaScript enabled.
  \item An internet connection.
  \item No account, no login, no installation.
\end{itemize}

\subsection{Opening the Dashboard}

\begin{enumerate}
  \item Open the browser and enter the address\\[2pt]
        \AppLink
  \item Wait until the chart appears (usually 2--5 seconds, Figure~\ref{fig:cloudLanding}). While the result files are being loaded, Streamlit shows a small running indicator.
  \item Bookmark the page for later use.
\end{enumerate}

\begin{figure}[H]
  \centering
  \Screenshot[0.9\linewidth]{Images/02Installation/CloudLanding.png}
  \caption{The dashboard opened from Streamlit Community Cloud}
  \label{fig:cloudLanding}
\end{figure}

\subsection{Waking the App Up}
\label{sec:instWake}

Streamlit Community Cloud puts apps to sleep when nobody has used them for a while \cite{StreamlitCloud:2026}. The first visitor after such a quiet period sees the page in Figure~\ref{fig:cloudSleeping} instead of the dashboard.

\begin{enumerate}
  \item Click the button \UI{Yes, get this app back up!}.
  \item Wait 30--60 seconds while the app starts. The page reloads by itself.
  \item The dashboard appears as usual. It stays awake for all users as long as it is being used.
\end{enumerate}

\begin{figure}[H]
  \centering
  \Screenshot[0.8\linewidth]{Images/02Installation/CloudSleeping.png}
  \caption{The page shown when the app is asleep}
  \label{fig:cloudSleeping}
\end{figure}

\begin{NoteBox}
Waking the app up does not change anything. The data and forecasts are the same as before; only the program is started again.
\end{NoteBox}

\section{Option B: Local Installation}
\label{sec:instLocal}

\subsection{Requirements}
\label{sec:instLocalReq}

\begin{table}[H]
\centering
\caption{Requirements for a local installation}
\label{tab:localReq}
\small
\begin{tabularx}{\linewidth}{@{}lX@{}}
\toprule
\textbf{Item} & \textbf{Requirement} \\
\midrule
Operating system & Windows 10/11, macOS 12 or newer, or a current Linux distribution \\
Python & Version 3.10 or newer (the dashboard uses language features introduced in 3.10) \\
Git & Any current version, to download the repository \\
Memory & 2\,GB of free \ac{ram} (the four result files are loaded completely) \\
Disk space & about 500\,MB for the repository and the Python packages \\
Network & Only for the installation; the dashboard itself runs offline \\
\bottomrule
\end{tabularx}
\end{table}

Check the Python and Git versions in a terminal (Terminal on macOS/Linux, PowerShell on Windows):

\begin{lstlisting}[style=shell]
python3 --version     # Windows: python --version
git --version
\end{lstlisting}

\subsection{Step 1 -- Download the Repository}

\begin{lstlisting}[style=shell]
git clone https://github.com/Hemanth-Jp/MBIDA_Master_Thesis.git
cd MBIDA_Master_Thesis
\end{lstlisting}

All following commands are run from this folder, the \emph{repository root}.

\begin{WarningBox}[Stay in the repository root]
Always start the dashboard from the repository root, not from inside \texttt{Code/app/}. Only then does Streamlit pick up the dashboard's settings file \texttt{.streamlit/config.toml}, which lives in the root folder -- exactly as on Streamlit Community Cloud.
\end{WarningBox}

\subsection{Step 2 -- Create a Virtual Environment}

A virtual environment keeps the dashboard's packages separate from everything else on the computer.

\begin{lstlisting}[style=shell]
python3 -m venv .venv

# activate it - macOS / Linux:
source .venv/bin/activate
# activate it - Windows (PowerShell):
.venv\Scripts\Activate.ps1
\end{lstlisting}

After activation, the prompt starts with \texttt{(.venv)}. If you prefer conda, the equivalent is \texttt{conda create -n m5\_dashboard python=3.10} followed by \texttt{conda activate m5\_dashboard}.

\subsection{Step 3 -- Install the Packages}

The dashboard needs only five packages. Their exact versions are listed in \\
 \texttt{Code/app/requirements.txt} 
 -- the same file Streamlit Community Cloud installs from (Section~\ref{sec:specSoftware}).

\begin{lstlisting}[style=shell]
pip install -r Code/app/requirements.txt
\end{lstlisting}

\begin{figure}[H]
  \centering
  \Screenshot[0.9\linewidth]{Images/02Installation/LocalPipInstall.png}
  \caption{Successful installation of the five packages}
  \label{fig:localPipInstall}
\end{figure}

\begin{WarningBox}[Package name]
The package that lets the chart follow the light/dark theme is called \texttt{st-theme}. A different package with the similar name \texttt{streamlit-theme} exists and does \emph{not} work. Always install from the requirements file instead of typing package names by hand.
\end{WarningBox}

\subsection{Step 4 -- Check the Result Files}

The dashboard only displays results; it needs the four result files in \\
 \texttt{Code/artifacts/dashboard\_data/}. They are part of the repository and are downloaded with it:

\begin{lstlisting}[style=shell]
ls Code/artifacts/dashboard_data     # Windows: dir Code\artifacts\dashboard_data
\end{lstlisting}

The listing must show \texttt{actuals.parquet}, \texttt{hierarchy.parquet}, \texttt{pred\_with\_gt.parquet} and \texttt{rec\_with\_gt.parquet}. If one of them is missing, regenerate them once as described in Section~\ref{sec:mtRegenerate}.

\subsection{Step 5 -- Start the Dashboard}

\begin{lstlisting}[style=shell]
streamlit run Code/app/streamlit_app.py
\end{lstlisting}

Streamlit prints the local address and opens it in the default browser (Figures~\ref{fig:localTerminal} and~\ref{fig:localBrowser}). On the very first start, Streamlit may ask for an e-mail address in the terminal; simply press \keys{Enter} to skip.

\begin{figure}[H]
  \centering
  \Screenshot[0.85\linewidth]{Images/02Installation/LocalTerminalLaunch.png}
  \caption{Terminal output after starting the dashboard locally}
  \label{fig:localTerminal}
\end{figure}

\begin{figure}[H]
  \centering
  \Screenshot[0.9\linewidth]{Images/02Installation/LocalBrowser.png}
  \caption{The local dashboard at \texttt{http://localhost:8501}}
  \label{fig:localBrowser}
\end{figure}

To stop the dashboard, press \keys{Ctrl+C} in the terminal. To start it again later, open a terminal in the repository root, activate the environment (Step~2) and repeat Step~5.

\section{Verifying the Installation}
\label{sec:instVerify}

The following check takes one minute and works for both options. If a step fails, go to the section named in the last column.

\begin{table}[H]
\centering
\caption{Installation check}
\label{tab:instCheck}
\small
\begin{tabularx}{\linewidth}{@{}clXl@{}}
\toprule
\textbf{\#} & \textbf{Action} & \textbf{Expected result} & \textbf{If not} \\
\midrule
1 & Open the dashboard & Title \UI{M5 Hierarchical Forecast Explorer}, no red error box & \ref{sec:tsParquet} \\
2 & Look at the node line & \UI{Selected node: Total (level = 1)} & \ref{sec:tsGeneral} \\
3 & Look at the chart & Black actual line, one dotted line, one solid line; three shaded regions & \ref{sec:tsNoRecLine} \\
4 & Set \UI{State} to \VAL{TX} & Chart changes, node line shows \VAL{Total/TX} & \ref{sec:tsGeneral} \\
5 & Set \UI{State} back to \VAL{(All)} & Grand total is shown again & -- \\
\bottomrule
\end{tabularx}
\end{table}

\begin{TipBox}
The keyboard keys in this manual (\keys{Ctrl+C}, \keys{Enter}) are the Windows/Linux names. On macOS, \keys{Ctrl+C} in the terminal is the same key combination (Control, not Command).
\end{TipBox}
